\subsection{The rest of a decoder block}

A common pre-normalized block is
\begin{align}
  U &= H+\operatorname{Attention}(\operatorname{Norm}(H)),\\
  H^{\mathrm{next}} &= U+\operatorname{MLP}(\operatorname{Norm}(U)).
\end{align}
Normalization controls the scale entering each learned sublayer. The residual additions
preserve a direct identity path, so a deep stack can modify a representation without
having to reconstruct it at every layer. These operations are linear in $BTD$ and are
usually small compared with the matrix multiplications.

The MLP acts independently at every token position. In its simplest form,
\begin{equation}
  \operatorname{MLP}(u)=W_2\,\sigma(W_1u),
  \qquad W_1\in\mathbb{R}^{F\times D},\quad
  W_2\in\mathbb{R}^{D\times F}.
\end{equation}
It contains approximately $2DF$ parameters and costs $4BTDF$ FLOPs. If $F=4D$, this is
$8D^2$ parameters and $16BTD^2$ FLOPs. The MLP often performs more arithmetic than
attention projections even though attention is the component that mixes information
across positions. Gated MLP variants use a third projection and therefore change these
constants, but not the basic token-wise structure.

Combining standard attention with a two-matrix MLP gives, per block,
\begin{align}
  \operatorname{parameters} &\approx 4D^2+2DF,\\
  \operatorname{forward\ FLOPs} &\approx 8BTD^2+4BT^2D+4BTDF.
\end{align}
With $F=4D$, these become approximately $12D^2$ parameters and
$24BTD^2+4BT^2D$ forward FLOPs. This decomposition is more informative than a single
model-wide number: increasing width mainly enlarges learned projections, while increasing
context length eventually exposes the quadratic attention term.
